% ============================================================================
%  Brouillon : ancienne sous-partie « Forme matricielle par blocs » du
%  chapitre 3, gardée comme matériau pour le chapitre 4. Non incluse.
% ============================================================================

\subsubsection{Forme matricielle par blocs}

Au terme constant $-\,m g a$ près, le lagrangien se réduit à l'énergie
cinétique, quadratique en $\dot q$ :
$T = \tfrac12\,\dot q^T\mat{M}(q)\,\dot q$. En généralisant l'exemple
$L = \tfrac12\,M(q)\,v^2 - V(q)$ du chapitre 1, les cinq équations
d'Euler-Lagrange-d'Alembert s'écrivent
\[
    \mat{M}(q)\,\ddot q + \vect{h}(q,\dot q) = \mat{A}(q)^T\vect{\lambda},
    \qquad
    \vect{h} = \dot{\mat{M}}\,\dot q
    - \tfrac12\,\nabla_q\bigl(\dot q^T\mat{M}\,\dot q\bigr)
\]
où $\vect{h}$ regroupe les termes quadratiques en $\dot q$. En dérivant
la contrainte $\mat{A}(q)\,\dot q = 0$ par rapport au temps,
\[
    \mat{A}(q)\,\ddot q + \dot{\mat{A}}\,\dot q = 0
\]

Pour expliciter ces termes, on sépare les coordonnées de translation
et d'orientation, $q = (r,\alpha)$ avec $r = (x,y)$ et
$\alpha = (\varphi,\theta,\psi)$, et on note $\mat{E}(\alpha)$ la
matrice de I.4.3 telle que $\vect{\omega} = \mat{E}(\alpha)\,\dot\alpha$.
On note aussi $\Id_2$ l'identité de taille $2$.

\paragraph{Matrice de masse}
Avec $T = \tfrac12\,m\,\dot r^T\dot r
+ \tfrac12\,\vect{\omega}^T\mat{I}_b\,\vect{\omega}$,
\[
    \mat{M} = \begin{pmatrix}
        m\,\Id_2 & \mat{0}\\
        \mat{0} & \mat{M}_\alpha
    \end{pmatrix},
    \qquad
    \mat{M}_\alpha = \mat{E}^T\mat{I}_b\,\mat{E}
\]
avec
\[
    \mat{M}_\alpha = \begin{pmatrix}
        I_1 & 0 & -I_1 s_\theta\\
        0 & I_2 c_\varphi^2 + I_3 s_\varphi^2
            & (I_2 - I_3)\,s_\varphi c_\varphi c_\theta\\
        -I_1 s_\theta & (I_2 - I_3)\,s_\varphi c_\varphi c_\theta
            & I_1 s_\theta^2 + (I_2 s_\varphi^2 + I_3 c_\varphi^2)\,
            c_\theta^2
    \end{pmatrix}
\]

\paragraph{Termes quadratiques}
Seul $\mat{M}_\alpha$ dépend de $q$, donc
$\vect{h} = (\vect{0},\vect{h}_\alpha)$. En pratique, on calcule
$\vect{h}_\alpha$ plus simplement à partir des équations d'Euler : un
couple $\vect{\tau}$ exprimé dans $\rep{b}$ a pour puissance
$\vect{\tau}\cdot\vect{\omega} = (\mat{E}^T\vect{\tau})\cdot\dot\alpha$,
donc les équations en $\alpha$ sont les équations d'Euler multipliées
par $\mat{E}^T$. Avec $\dot{\vect{\omega}} = \mat{E}\,\ddot\alpha
+ \dot{\mat{E}}\,\dot\alpha$,
\[
    \vect{h}_\alpha = \mat{E}^T\bigl(\mat{I}_b\,\dot{\mat{E}}\,\dot\alpha
    + \vect{\omega}\wedge(\mat{I}_b\,\vect{\omega})\bigr),
    \qquad
    \dot{\mat{E}}\,\dot\alpha = \begin{pmatrix}
        -\,\dot\theta\dot\psi\,c_\theta\\
        -\,\dot\varphi\dot\theta\,s_\varphi
            + \dot\psi\,(\dot\varphi\,c_\varphi c_\theta
            - \dot\theta\,s_\varphi s_\theta)\\
        -\,\dot\varphi\dot\theta\,c_\varphi
            - \dot\psi\,(\dot\varphi\,s_\varphi c_\theta
            + \dot\theta\,c_\varphi s_\theta)
    \end{pmatrix}
\]

\paragraph{Contrainte}
La matrice des contraintes se découpe en
$\mat{A} = (\Id_2\ \ \mat{A}_\alpha)$, et seul $\mat{A}_\alpha$ dépend
de $q$, donc $\dot{\mat{A}}\,\dot q = \dot{\mat{A}}_\alpha\,\dot\alpha$,
avec
\[
    \mat{A}_\alpha = a\begin{pmatrix}
        -s_\psi c_\theta & -c_\psi & 0\\
        c_\psi c_\theta & -s_\psi & 0
    \end{pmatrix},
    \qquad
    \dot{\mat{A}}_\alpha\,\dot\alpha = a\begin{pmatrix}
        -\,\dot\varphi\dot\psi\,c_\psi c_\theta
            + \dot\varphi\dot\theta\,s_\psi s_\theta
            + \dot\theta\dot\psi\,s_\psi\\
        -\,\dot\varphi\dot\psi\,s_\psi c_\theta
            - \dot\varphi\dot\theta\,c_\psi s_\theta
            - \dot\theta\dot\psi\,c_\psi
    \end{pmatrix}
\]

\begin{proposition}[Forme par blocs]
    Pour des conditions initiales qui vérifient la contrainte, le
    système du théorème précédent équivaut à
    \[
        \begin{pmatrix}
            \mat{M}(q) & -\,\mat{A}(q)^T\\
            \mat{A}(q) & \mat{0}
        \end{pmatrix}
        \begin{pmatrix}
            \ddot q\\ \vect{\lambda}
        \end{pmatrix}
        = \begin{pmatrix}
            -\,\vect{h}(q,\dot q)\\ -\,\dot{\mat{A}}\,\dot q
        \end{pmatrix}
    \]
    et la matrice du système est inversible tant que
    $\cos\theta \neq 0$.
\end{proposition}

\begin{demonstration}
    La première ligne de blocs regroupe les équations
    d'Euler-Lagrange-d'Alembert, la seconde est la contrainte dérivée,
    qui équivaut à la contrainte si elle est vérifiée à l'instant
    initial. Pour l'inversibilité, $\mat{M}_\alpha = \mat{E}^T\mat{I}_b\,
    \mat{E}$ est symétrique définie positive car $\mat{I}_b$ l'est et
    $\det\mat{E} = \cos\theta \neq 0$, donc $\mat{M}$ aussi. Soit
    $(u,\vect{\nu})\in\R^5\times\R^2$ dans le noyau, c'est-à-dire
    $\mat{M}\,u = \mat{A}^T\vect{\nu}$ et $\mat{A}\,u = 0$. Alors
    \[
        u^T\mat{M}\,u = u^T\mat{A}^T\vect{\nu}
        = (\mat{A}\,u)^T\vect{\nu} = 0
    \]
    donc $u = 0$, puis $\mat{A}^T\vect{\nu} = \vect{0}$, donc
    $\vect{\nu} = \vect{0}$ car $\rang\mat{A} = 2$. La matrice, carrée
    et de noyau nul, est inversible.
\end{demonstration}

Ainsi, tant que $\cos\theta \neq 0$, $(q,\dot q)$ détermine
$(\ddot q,\vect{\lambda})$ à chaque instant : c'est la forme utilisée
pour la simulation numérique.

\begin{corollaire}[Système réduit]
    En éliminant $\vect{\lambda} = m\,\ddot r$ et
    $\ddot r = -(\mat{A}_\alpha\,\ddot\alpha
    + \dot{\mat{A}}_\alpha\,\dot\alpha)$, l'orientation vérifie
    \[
        \bigl(\mat{M}_\alpha + m\,\mat{A}_\alpha^T\mat{A}_\alpha\bigr)\,
        \ddot\alpha
        = -\,\vect{h}_\alpha - m\,\mat{A}_\alpha^T\dot{\mat{A}}_\alpha\,
        \dot\alpha
    \]
    puis $\dot r = -\,\mat{A}_\alpha\,\dot\alpha$ et
    $\vect{\lambda} = -\,m\,(\mat{A}_\alpha\,\ddot\alpha
    + \dot{\mat{A}}_\alpha\,\dot\alpha)$.
\end{corollaire}

La matrice $\mat{M}_\alpha + m\,\mat{A}_\alpha^T\mat{A}_\alpha$ est la
matrice d'inertie autour du point de contact, écrite dans les
coordonnées $\alpha$. En effet, $\vect{v}_G = a\,\vect{\Omega}\wedge
\vect{e}_z$ donne $m\,\|\vect{v}_G\|^2 = m\,a^2\,\vect{\Omega}^T
(\Id - \vect{e}_z\vect{e}_z^T)\,\vect{\Omega}$, donc, pour un mouvement
qui respecte la contrainte,
\[
    T = \tfrac12\,\dot\alpha^T\bigl(\mat{M}_\alpha
    + m\,\mat{A}_\alpha^T\mat{A}_\alpha\bigr)\,\dot\alpha
    = \tfrac12\,\vect{\Omega}^T\mat{I}_P\,\vect{\Omega},
    \qquad
    \mat{I}_P = \mat{I}_G
    + m\,a^2\,\bigl(\Id - \vect{e}_z\vect{e}_z^T\bigr)
\]
ce qui est le théorème de Huygens-Steiner du chapitre 2 avec
$\overrightarrow{PG} = a\,\vect{e}_z$.
